%% ── Chapter 1: Introduction ──────────────────────────────────
\chapter{Introduction}
\label{ch:introduction}

\section{Motivation and Context}

The digital transformation of critical infrastructure has created an
unprecedented attack surface. Adversaries range from opportunistic
script-kiddies to nation-state actors deploying advanced persistent threats
(APTs) that combine multiple attack vectors, adapt in real-time to defensive
measures, and persist undetected for months.
Traditional cybersecurity relies on \emph{static} defenses:
signature-based IDS, rule-based firewalls, and periodic vulnerability scanning.
These systems are fundamentally reactive — they can only defend against
attacks that have been observed before.

Meanwhile, the rise of generative artificial intelligence has given
adversaries powerful new tools.
Large language models can craft convincing spear-phishing emails at scale.
Generative Adversarial Networks can synthesize malicious network flows
indistinguishable from benign traffic.
Diffusion models can perturb malware to evade static analysis.
The asymmetry is stark: attackers leverage AI to \emph{innovate},
while defenders rely on catalogues of \emph{known} signatures.

This thesis addresses the fundamental challenge:
\begin{quote}
\textit{Can we build AI systems that \textbf{continuously simulate}
realistic attacks and \textbf{adaptively evolve} defenses in a closed
co-evolutionary loop, ultimately achieving autonomous, self-healing security?}
\end{quote}

\section{Research Questions}

This thesis investigates the following research questions:

\begin{description}
  \item[RQ1] Can a defense-conditioned generative model produce attack traffic
    that adapts to deployed countermeasures, better mimicking real APT behaviour
    than unconditional generators?

  \item[RQ2] Can a reinforcement learning agent trained on co-evolutionarily
    generated attack data achieve superior detection and response performance
    compared to classical IDS baselines?

  \item[RQ3] Does a temporal graph of attack-defense interactions enable
    statistically significant proactive threat anticipation?

  \item[RQ4] Can a diffusion-based adversarial perturbation engine produce
    attack samples that simultaneously evade both classifier-based and
    anomaly-based defenses?

  \item[RQ5] Is fully autonomous incident response (zero-touch self-healing)
    achievable within a co-evolutionary AI security framework?
\end{description}

\section{Contributions}

This thesis makes the following novel scientific contributions:

\begin{enumerate}
  \item \textbf{The \caesar{} framework}: the first unified co-evolutionary
    architecture for simultaneous attack simulation and adaptive defense
    (\cref{ch:framework}).

  \item \textbf{\tagan{}}: a defense-conditioned GAN whose generator
    $G(z, \mathbf{d}) \to (\text{attack\_type}, \text{intensity})$ conditions
    on the real-time defense embedding $\mathbf{d} \in \R^8$,
    producing attacks that adapt to what the defender is currently deploying
    (\cref{ch:tagan}).

  \item \textbf{\adpn{}}: a Dueling Double-DQN agent with a co-evolutionary
    reward signal $r_{\mathrm{def}} = R_{\mathrm{env}} \cdot (1 - \alpha \cdot
    r_a) + \beta \cdot H$ linking defense quality to both environment feedback
    and attacker performance (\cref{ch:adpn}).

  \item \textbf{\tagraph{}}: a directed temporal knowledge graph of attack-defense
    co-occurrences supporting Markov-chain attack prediction and
    confidence-gated proactive countermeasure deployment (\cref{ch:tag}).

  \item \textbf{\mgae{}}: a Manifold-Guided Adversarial Engine using a cosine
    denoising diffusion process to project attack samples onto the learned
    benign traffic manifold, achieving simultaneous evasion of classifier-based
    and anomaly-based IDS (\cref{ch:mgae}).

  \item \textbf{\arl{}}: an Autonomous Remediation Loop implementing a
    five-state healing state machine with countermeasure rotation, healing
    confirmation, and severity-based escalation (\cref{ch:selfhealing}).
\end{enumerate}

\section{Thesis Structure}

The remainder of this thesis is organised as follows.
\Cref{ch:background} provides mathematical background on GANs, reinforcement
learning, and diffusion models.
\Cref{ch:related} surveys related work.
\Cref{ch:framework} presents the \caesar{} framework architecture.
Chapters~\ref{ch:tagan}--\ref{ch:selfhealing} detail each novel component.
\Cref{ch:experiments} presents experimental results.
\Cref{ch:discussion} discusses implications and limitations.
\Cref{ch:conclusion} concludes with future directions.

\section{Publications and Artifacts}

The source code for \caesar{} is available at:
\begin{center}
  \url{https://github.com/gdtsitlauri/caesar-framework}
\end{center}

Portions of this work have been submitted to / accepted at:
\begin{itemize}
  \item \textit{[Conference/Journal Name]}, [Year] — \textit{[Paper Title]}
\end{itemize}
